\begin{abstract}
Why do organizational transformations fail at rates approaching 90\%? Why can AI suddenly replace knowledge workers who spent decades developing expertise? Because we spent centuries training humans to be biological AI systems—and now the silicon versions have arrived to collect their inheritance. I spent thirty years inside this process, building trading platforms, optimizing business systems, implementing the frameworks I now diagnose. This paper is confession literature from someone who helped prepare the feast.

Using a thermodynamic analysis, we show that knowledge requires continuous energy investment to maintain its complexity. Without this investment, expertise degrades to procedures, then to rituals, and finally to patterns simple enough for algorithms to replicate. We trace this degradation from ancient Greek academies, through medieval universities, to today's micro-credentials. The 1999 Bologna Process industrialized this collapse, fragmenting knowledge into tradeable credits while eliminating the comprehensive foundations that once preceded specialization---a transformation whose architects now acknowledge produced graduates unqualified for professional practice.

Our framework reveals how education and management systems systematically transformed human cognition from multidimensional ``spheres''---capable of navigating complexity through cross-domain synthesis---into one-dimensional ``vectors'' optimized for procedural efficiency. The thermodynamic equation, Cognitive Sovereignty = (Energy Invested/Time) $\times$ Resistance to Extraction, explains both the 42\% AI project abandonment rate and why leading German technical universities are reviving the Diplom-Ingenieur they sacrificed to European standardization, while 6-week bootcamps gleefully produce GPT-5.x fodder: vectors can optimize, but only spheres can adapt and connect the dots across different domains.

The implications are stark. We haven't reached ``peak human'' and been surpassed by superior AI. We've depleted human cognitive development until simple pattern-matching systems can replicate our degraded outputs. The thermodynamics are simple: invest energy in cognitive development or watch expertise degrade until algorithms can replicate what's left. There is no pause button. Entropy is already working.

Resistance to the 2\textsuperscript{nd} law of thermodynamics is futile---Physics doesn't negotiate.
\end{abstract}